\documentclass[12pt]{article}

% Template SBC fornecido no projeto
\usepackage{sbc-template}

% Pacotes já utilizados pelo modelo e suporte às equações
\usepackage{graphicx}
\usepackage{url}
\usepackage[utf8]{inputenc}
\usepackage[T1]{fontenc}
\usepackage[brazil]{babel}

\title{Classificação de variedades de feijão por uma rede Perceptron multicamadas em NumPy}

\author{
  Patrick Correa Moraes\inst{1}
}

\address{
  Universidade Estadual do Centro-Oeste (UNICENTRO)\\
  Cidade -- Paraná -- Brasil
  \email{patrick.pitanga2017@gmail.com}
}

\begin{document}

\maketitle

\begin{abstract}
This paper describes a NumPy implementation of a multilayer perceptron for classifying seven dry bean varieties from sixteen numerical descriptors. The Dry Bean dataset was divided into stratified training, validation, and test sets. The network uses two hidden layers, ReLU activations, a softmax output, categorical cross-entropy, and manually implemented backpropagation with Adam updates. In the recorded run, the model correctly classified 1,881 of 2,042 test samples, corresponding to an accuracy of 92.12\%. The experiment emphasizes transparent implementation and a reproducible evaluation protocol.
\end{abstract}

\begin{resumo}
Este artigo apresenta uma implementação em NumPy de uma rede Perceptron multicamadas para classificar sete variedades de feijão a partir de dezesseis descritores numéricos. O Dry Bean Dataset foi dividido de forma estratificada em treino, validação e teste. A rede possui duas camadas intermediárias, ativações ReLU, saída softmax, entropia cruzada e retropropagação implementada manualmente, com atualização por Adam. Na execução registrada, foram classificadas corretamente 1.881 das 2.042 amostras de teste, o que corresponde a 92,12\% de acurácia. O trabalho prioriza a transparência do algoritmo e a avaliação reproduzível.
\end{resumo}

\section{Introdução}

A classificação automática de sementes pode apoiar tarefas de inspeção, padronização e análise de qualidade na cadeia agrícola. A variedade de um grão está associada a propriedades físicas observáveis, mas variedades distintas podem apresentar dimensões e formas semelhantes. Por isso, atribuir uma classe somente por inspeção visual pode ser uma tarefa trabalhosa e sujeita à variação entre observadores. Um classificador computacional oferece uma maneira de relacionar medidas numéricas do grão com uma variedade previamente definida, desde que existam exemplos rotulados que permitam ajustar e avaliar o modelo.

Este trabalho aborda a classificação supervisionada de sete variedades de feijão seco conhecidas como Barbunya, Bombay, Cali, Dermason, Horoz, Seker e Sira. A tarefa utiliza descritores geométricos já extraídos de imagens, em vez de processar as imagens diretamente. Cada amostra é, portanto, representada por dezesseis números e por um rótulo de classe. O objetivo é aprender uma função que receba os descritores e estime a variedade correspondente, mantendo a coluna de classe fora das entradas da rede. O conjunto de dados foi originalmente construído a partir de imagens de grãos e contém medições morfológicas de 13.611 amostras.

A solução implementa uma rede neural Perceptron multicamadas, ou MLP, com operações matriciais feitas manualmente por meio de NumPy. A implementação não delega o treinamento a uma biblioteca de aprendizagem profunda nem usa uma rotina pronta de gradientes. O código calcula a propagação para frente, a entropia cruzada, a retropropagação e as atualizações dos pesos e biases. Bibliotecas auxiliares são utilizadas para particionar e normalizar os dados e para calcular métricas, preservando o foco didático no funcionamento interno da rede.

A arquitetura usada na execução apresentada tem dezesseis neurônios de entrada, duas camadas intermediárias com dezesseis e oito neurônios, respectivamente, e sete neurônios de saída. As camadas intermediárias empregam ReLU e a camada final emprega softmax, de modo que as saídas possam ser interpretadas como probabilidades das classes. O treinamento é feito em lotes, com embaralhamento dos exemplos a cada época e atualização por Adam. Uma parte dos dados é reservada à validação para selecionar a melhor época, enquanto o teste é mantido separado até a avaliação final.

A contribuição principal do projeto é apresentar um caminho verificável entre os dados e a previsão. O artigo acompanha os módulos e métodos que carregam os dados, montam as camadas, calculam as probabilidades, ajustam os parâmetros e avaliam o modelo. A avaliação não se limita a apresentar a acurácia: inclui também precisão, recall, F1-score e análise da matriz de confusão, permitindo identificar classes bem reconhecidas e pares que geram erros. A tabela de experimentos é mantida sem valores para permitir o registro posterior de execuções com diferentes arquiteturas e taxas de aprendizado.

O restante do artigo organiza essa explicação em quatro partes. A seção sobre a implementação acompanha a organização da MLP e o percurso de uma amostra pelo código, do processamento inicial à atualização dos pesos. A metodologia documenta os dados, a separação dos conjuntos e as decisões de implementação. Em seguida, os resultados da execução registrada são analisados junto à tabela reservada para novos experimentos. A conclusão discute os acertos, as limitações e a utilidade da implementação para o estudo de redes neurais. A estrutura e a apresentação deste manuscrito seguem o modelo de artigos da Sociedade Brasileira de Computação e suas instruções para autores \cite{nedel2017sbc}.

\section{Funcionamento da MLP implementada}

\subsection{Organização interna: neurônios, conexões e camadas}

A rede implementada recebe uma linha de atributos por feijão. Essa linha contém dezesseis valores numéricos, como área, perímetro, comprimento dos eixos e medidas de forma. O programa não interpreta cada valor como uma classe ou como uma resposta independente. Em vez disso, apresenta os dezesseis valores juntos à rede, que os processa em sequência até produzir uma distribuição de probabilidades para as sete variedades. O nome da variedade verdadeira permanece separado das entradas e só é consultado durante o treinamento e a avaliação.

Na implementação, um neurônio é representado pelos valores armazenados nas matrizes de pesos e nos vetores de biases. Cada conexão entre duas camadas possui um peso, que é um número ajustado durante o treinamento. Esse valor determina como uma informação recebida de um neurônio influencia o cálculo do neurônio seguinte. Cada neurônio treinável também possui um bias próprio. O bias é somado ao resultado das conexões e permite que o neurônio ajuste seu ponto de resposta mesmo quando os valores recebidos são próximos de zero. Tanto pesos quanto biases são parâmetros que começam com valores iniciais e mudam ao longo do treinamento.

As camadas intermediárias são compostas por grupos de neurônios que transformam a representação recebida. A arquitetura utilizada no experimento possui uma primeira camada intermediária com dezesseis neurônios e uma segunda com oito. As ligações são densas: cada valor da entrada participa do cálculo de todos os neurônios da primeira camada; cada saída da primeira camada participa do cálculo de todos os neurônios da segunda; e cada neurônio da segunda camada se conecta às sete unidades de saída. Essa organização permite que a rede combine informações de diferentes características antes de decidir qual variedade é mais provável.

A classe \texttt{PerceptronMulticamadas}, definida em \texttt{perceptron\_multicamadas.py}, cria essa arquitetura a partir dos tamanhos informados na configuração. Ela guarda os pesos em uma lista de matrizes e os biases em uma lista de vetores. Na arquitetura do experimento, a primeira matriz de pesos tem dezesseis linhas e dezesseis colunas, a seguinte tem dezesseis linhas e oito colunas, e a matriz final tem oito linhas e sete colunas. As dimensões indicam quantos valores saem de uma camada e quantos neurônios existem na camada seguinte. O uso de matrizes permite fazer o cálculo de uma camada inteira com operações do NumPy, sem escrever um laço Python para cada conexão.

A lista de camadas intermediárias é variável. Se o estudante altera essa lista no arquivo \texttt{config.py}, o construtor da MLP cria as matrizes e os vetores correspondentes à nova sequência. Por exemplo, uma lista com uma única quantidade cria apenas uma camada intermediária; uma lista com três quantidades cria três. O tamanho da entrada continua sendo validado contra o número de características do arquivo e a quantidade de saídas continua sendo comparada ao número de classes. O programa interrompe a execução com uma mensagem explicativa se essas dimensões não forem compatíveis, evitando que a rede altere a arquitetura por conta própria.

\subsection{Inicialização e propagação para frente}

Antes do treinamento, o construtor inicializa os parâmetros. Quando uma camada intermediária utiliza ReLU, seus pesos são sorteados, o que é adequado para essa ativação. Para a camada de saída, a implementação usa uma escala de inicialização. Os biases começam em zero. O sorteio é realizado por um gerador do NumPy criado com a seed definida pelo usuário; assim, a mesma arquitetura e a mesma seed produzem a mesma inicialização. Essa repetição ajuda a comparar execuções porque evita que uma mudança de resultado venha apenas de um sorteio inicial diferente.

A propagação para frente é executada pelo método \texttt{propaga\_para\_frente}. Ele recebe uma matriz NumPy com uma amostra em cada linha e verifica se todas as linhas têm dezesseis características e se todos os valores são finitos. Em seguida, percorre as matrizes de pesos e os vetores de biases, camada por camada. Para cada camada, o código combina os valores atuais com os pesos e soma os biases; o trecho central dessa operação é \texttt{peso + bias}. O resultado intermediário é enviado à ativação correspondente e passa a ser a entrada da próxima camada.

Quando o método é chamado durante o treinamento, ele também devolve um cache. Esse cache guarda a entrada original, as ativações produzidas em cada camada e os valores calculados antes das ativações. O cache não representa uma segunda previsão nem modifica os dados. Ele conserva informações da passagem atual para que o método de retropropagação consiga calcular como cada parâmetro participou da saída. Quando a rede está apenas classificando novas amostras, o cache não é necessário e o método devolve somente as probabilidades finais.

Nas camadas intermediárias, a configuração padrão utiliza ReLU. No código, essa ativação substitui por zero os valores negativos e mantém os positivos. Esse comportamento permite que a rede combine as camadas de maneiras mais variadas do que uma simples transformação linear. O projeto também implementa sigmoid como opção para as camadas intermediárias. A configuração informa qual ativação deve ser usada, e a rede chama o método correspondente tanto durante a previsão quanto durante a retropropagação. No experimento relatado, a opção ativa foi ReLU.

Na saída, a MLP usa softmax. Antes de calcular as exponenciais, o método subtrai de cada linha o maior valor calculado para aquela amostra. Esse passo protege a operação contra valores numéricos excessivamente grandes. Depois, normaliza os resultados para que cada linha se transforme em sete probabilidades cuja soma fica próxima de um. Cada posição está associada a uma classe do mapeamento salvo no modelo. O programa escolhe como previsão a posição de maior probabilidade, mas conserva todas as probabilidades para calcular a perda e permitir uma inspeção mais detalhada da resposta.

\subsection{Perda e retropropagação no código}

As classes textuais são convertidas pelo pré-processamento em índices inteiros de zero a seis. A rede recebe esses índices como respostas esperadas, sem construir uma matriz permanente com sete posições para cada rótulo. A função \texttt{perda} seleciona, para cada amostra, a probabilidade atribuída à classe verdadeira e calcula a entropia cruzada média do lote. Uma previsão que atribui baixa probabilidade à classe correta contribui com uma perda maior. O cálculo limita a probabilidade a um mínimo positivo antes do logaritmo, para evitar um resultado infinito quando a probabilidade numérica for zero. Se a regularização L2 estiver ativa, a função também inclui a penalização dos pesos; ela não inclui os biases.

A retropropagação está concentrada no método \texttt{backpropagation}. Ele recebe os rótulos inteiros e o cache criado pela propagação para frente. Na saída, começa com uma cópia das probabilidades previstas. Para cada linha, reduz em um o valor que corresponde à classe verdadeira, usando a instrução \texttt{delta[np.arange(batch\_size), y\_true] -= 1.0}. O resultado equivale ao erro inicial usado na combinação entre softmax e entropia cruzada, sem materializar uma matriz de rótulos one-hot. Essa escolha reduz alocações e mantém os rótulos no formato simples de índices inteiros.

Depois de preparar esse erro na saída, o método percorre as camadas em ordem inversa. Em cada uma, combina as ativações guardadas no cache com o erro que chegou à camada e calcula uma matriz de gradientes para os pesos. Também calcula um gradiente para cada bias, somando a contribuição dos exemplos do lote. As contribuições são divididas pelo tamanho do lote, de modo que os gradientes representem a média das amostras processadas naquela atualização. Se L2 estiver ativada, a função acrescenta a contribuição correspondente aos pesos. Os biases continuam sem penalização.

Para continuar em direção às camadas anteriores, o erro é propagado pelas conexões da camada seguinte e ajustado pela derivada da ativação usada na camada atual. No caso da ReLU, a implementação cria uma máscara a partir dos valores guardados antes da ativação: os neurônios que receberam valor positivo transmitem gradiente, enquanto os que ficaram inativos não o transmitem naquele passo. Se a opção sigmoid estiver selecionada, o método usa os próprios valores de ativação para calcular sua derivada. Como os dados intermediários foram guardados durante a propagação para frente, não é necessário executar novamente a rede para produzir essas informações.

A retropropagação, portanto, é o trecho que calcula gradientes para todos os pesos e biases; ela não atualiza sozinha os parâmetros. Essa separação aparece no código: \texttt{backpropagation} devolve listas de gradientes, enquanto \texttt{atualiza\_parametros} recebe essas listas e altera os parâmetros. A combinação da regra da cadeia com a organização em camadas segue a abordagem clássica de treinamento de redes multicamadas \cite{rumelhart1986}, mas nesta implementação os cálculos são expressos diretamente em operações NumPy, sem autograd, Keras ou TensorFlow.

\subsection{Atualização dos parâmetros, épocas e lotes}

O projeto oferece dois métodos de atualização. Com SGD, cada parâmetro é alterado diretamente na direção indicada pelo gradiente e pela taxa de aprendizado. Com Adam, o método \texttt{atualiza\_parametros} mantém, para cada peso e bias, médias móveis do gradiente e dos seus quadrados. Também corrige o efeito de início dessas médias e usa o resultado para definir o tamanho de cada atualização. A implementação usa os valores padrão de Adam para os coeficientes das médias e a proteção numérica, enquanto a taxa de aprendizado vem do arquivo de configuração. No experimento descrito, a opção foi Adam \cite{kingma2015}. Backpropagation calcula os gradientes; Adam usa esses gradientes para atualizar os parâmetros, portanto são etapas diferentes.

O tamanho do lote controla quantos exemplos contribuem para uma atualização. O método \texttt{treina\_lote} reúne a propagação para frente, o cálculo da perda, a retropropagação e a atualização dos parâmetros. O módulo \texttt{treinamento.py} chama esse método repetidamente. No início de cada época, ele embaralha os índices das amostras de treino e separa os índices em lotes de até 32 exemplos. Se a quantidade restante no último lote for menor, esse lote também é processado. Assim, todas as amostras de treino participam de cada época, ainda que os lotes não tenham exatamente o mesmo tamanho.

Uma época termina quando todos os lotes de treino foram processados. Em seguida, o treinamento calcula perda e acurácia nos conjuntos de treino e validação. Essas medições são feitas sem atualizar os parâmetros. O histórico guarda os valores para cada época e é usado tanto no acompanhamento do terminal quanto na geração posterior dos gráficos. A acurácia mostra a proporção de previsões corretas; a perda também considera se a rede atribuiu probabilidade alta ou baixa à classe verdadeira. Por isso, as duas medidas podem evoluir de maneiras diferentes ao longo do ajuste.

O módulo também implementa parada antecipada. Depois de cada época, compara a perda de validação com o melhor valor registrado. Se houver melhora superior ao limite mínimo configurado, salva uma cópia dos pesos e biases e reinicia a contagem de épocas sem melhora. Caso contrário, aumenta essa contagem. Quando o limite de paciência é atingido, o laço encerra mesmo que ainda existam épocas disponíveis. No fim, o programa restaura a cópia que selecionou pela validação. Essa restauração garante que o teste avalie a versão escolhida durante o desenvolvimento, em vez de simplesmente usar os parâmetros da última época.

Na execução descrita, o limite foi de cem épocas e a paciência foi de quinze épocas, com uma melhora mínima de 0,0001. O treinamento encerrou após 59 épocas e restaurou os parâmetros da época 44, que teve a menor perda de validação no histórico registrado. A seed é usada para a inicialização da MLP, para as divisões dos dados e para o embaralhamento das amostras. Mantendo iguais o arquivo, a configuração e as bibliotecas, o programa repete o mesmo fluxo aleatório e facilita a comparação entre execuções.

\subsection{Da avaliação à previsão de uma amostra nova}

O conjunto de teste não participa da retropropagação nem da escolha da época. Depois que o treinamento restaura o melhor estado, o módulo \texttt{avaliacao.py} recebe as previsões do conjunto de teste e usa funções auxiliares do scikit-learn para calcular acurácia, precisão, recall, F1-score, relatório de classificação e matriz de confusão. Essas funções resumem as previsões; não atualizam a rede. A matriz registra quantas amostras de cada classe real foram atribuídas a cada classe prevista, permitindo identificar confusões específicas que uma média geral não mostra.

A rede não precisa ser treinada novamente para classificar uma amostra depois que o modelo foi salvo. O método \texttt{salva} grava os pesos e biases em um arquivo NumPy compactado e grava em JSON os dados necessários para reconstruir a arquitetura e interpretar as saídas. Entre esses dados estão a lista de características, os nomes das classes, o mapeamento dos rótulos, as opções de ativação, a seed e a configuração usada. Quando a normalização está ativa, os parâmetros aprendidos pelo scaler também são incluídos nos metadados.

Para usar o modelo em uma amostra nova, o usuário fornece os mesmos dezesseis atributos, na ordem identificada pelos nomes das colunas. O método \texttt{carrega} reconstrói a rede e lê os parâmetros salvos. O método \texttt{preve\_dados\_brutos} aplica aos novos valores a mesma padronização aprendida no conjunto de treino, calcula as probabilidades e converte a posição mais provável de volta ao nome da variedade. O resultado é o tipo previsto junto das probabilidades das sete classes. A previsão não altera pesos e não executa backpropagation: ela apenas repete a propagação para frente com o modelo já treinado.

\section{Metodologia}

\subsection{Descrição e inspeção dos dados}

O experimento utiliza o Dry Bean Dataset disponibilizado pelo UCI Machine Learning Repository. O conjunto reúne 13.611 grãos e dezesseis características numéricas derivadas de medidas geométricas de regiões segmentadas em imagens. O alvo é a variedade declarada na coluna \texttt{Class}, com sete valores: Seker, Barbunya, Bombay, Cali, Dermason, Horoz e Sira \cite{ucibean}. A descrição do conjunto também apresenta a construção de um sistema de visão e a extração dos atributos morfológicos, tendo como estudo introdutório o trabalho de Koklu e Ozkan \cite{koklu2020}.

As características representam tamanho, proporções e forma. A área e a área convexa medem a extensão do grão e do seu envoltório convexo; o perímetro registra a extensão do contorno; os comprimentos dos eixos maior e menor resumem suas dimensões e orientação; a relação entre os eixos e a excentricidade descrevem alongamento. O diâmetro equivalente relaciona a área a um círculo de mesma área; extensão e solidez descrevem a ocupação e a convexidade; arredondamento e compactação expressam propriedades de forma; e quatro fatores de forma completam a descrição. No arquivo, esses atributos aparecem como \texttt{Area}, \texttt{Perimeter}, \texttt{MajorAxisLength}, \texttt{MinorAxisLength}, \texttt{AspectRation}, \texttt{Eccentricity}, \texttt{ConvexArea}, \texttt{EquivDiameter}, \texttt{Extent}, \texttt{Solidity}, \texttt{roundness}, \texttt{Compactness}, \texttt{ShapeFactor1}, \texttt{ShapeFactor2}, \texttt{ShapeFactor3} e \texttt{ShapeFactor4}. A grafia \texttt{AspectRation} é mantida como encontrada no arquivo do dataset.

O arquivo XLSX original é carregado pelo módulo de leitura, que procura as características pelo nome e não depende da ordem em que as colunas aparecem. O carregador verifica o formato, a existência da coluna-alvo, a presença das dezesseis características, os valores ausentes, os valores numéricos e a quantidade de classes. Essa validação impede, por exemplo, que a coluna \texttt{Class} entre por engano como uma entrada numérica ou que um arquivo incompatível seja processado silenciosamente. Embora as imagens tenham sido utilizadas na origem para obter as medidas, a implementação deste artigo recebe somente a tabela de atributos e não executa segmentação ou extração visual.

\subsection{Particionamento e normalização}

A divisão foi configurada em 70\% para treino, 15\% para validação e 15\% para teste, com estratificação e semente aleatória 42. Na execução registrada, isso produziu 9.527 exemplos de treinamento, 2.042 de validação e 2.042 de teste. A estratificação busca manter proporções de classe semelhantes nos três subconjuntos, reduzindo a possibilidade de uma variedade pouco frequente ficar concentrada em apenas uma partição. A semente fixa a aleatoriedade da divisão, do embaralhamento por época e da inicialização, permitindo reproduzir o mesmo protocolo enquanto os dados e as configurações permanecerem iguais.

As características possuem unidades e escalas diferentes. Por exemplo, áreas e perímetros podem ter magnitudes muito superiores às de razões geométricas. Quando uma variável tem escala muito maior, ela pode dominar as atualizações numéricas e tornar o ajuste mais difícil. Por isso, a configuração usa o \texttt{StandardScaler}, que centraliza cada característica pela média do treino e ajusta sua escala pelo desvio padrão calculado nesse mesmo conjunto. O scaler é ajustado exclusivamente com os exemplos de treino; seus parâmetros aprendidos são aplicados sem novo ajuste aos conjuntos de validação e teste \cite{pedregosa2011}. Esse procedimento evita que informações estatísticas das partições de validação ou teste influenciem o pré-processamento. O scaler também é salvo junto ao modelo para transformar novas amostras da mesma maneira.

O pré-processamento também codifica cada classe textual em um índice inteiro estável. No experimento, a ordenação dos nomes foi usada para criar o mapeamento entre as sete variedades e os índices de saída de zero a seis. Esse mapeamento é guardado nos metadados do modelo, junto à lista ordenada das características. Assim, uma previsão numérica pode ser convertida de volta para o nome da variedade e uma nova amostra pode ser reordenada corretamente antes de entrar na rede.

\subsection{Arquitetura e inicialização}

A arquitetura experimental utiliza dezesseis entradas, duas camadas intermediárias com dezesseis e oito neurônios, e sete saídas. As ligações entre camadas adjacentes são completas, como ocorre na classe \texttt{PerceptronMulticamadas}. A configuração permite trocar a quantidade de camadas e de neurônios sem editar a implementação central; o construtor recria as matrizes de pesos e vetores de biases a partir da lista configurada. O código inicializa os pesos das camadas ReLU com a estratégia He e os da saída com a estratégia Xavier. Os biases começam em zero, e um gerador aleatório controlado pela seed produz os parâmetros iniciais.

Durante a execução, o método \texttt{propaga\_para\_frente} percorre todas as camadas e guarda os valores intermediários necessários ao treinamento. \texttt{perda} compara as probabilidades com os índices de classe verdadeiros. \texttt{backpropagation} percorre a rede no sentido inverso e produz gradientes para pesos e biases. Por fim, \texttt{atualiza\_parametros} aplica o otimizador escolhido; no experimento registrado, foi Adam. O método \texttt{treina\_lote} reúne essas operações para uma atualização usando um lote de amostras.

\subsection{Laço de treinamento e prevenção de vazamento}

O treinamento é coordenado por um laço que, em cada época, cria uma permutação dos índices de treino. Essa permutação divide os dados em lotes de 32 exemplos. Para cada lote, o modelo calcula as probabilidades, mede a perda, executa backpropagation para calcular os gradientes e atualiza seus parâmetros. A taxa de aprendizado foi definida como 0,001 e o limite máximo foi de cem épocas. A perda usada é entropia cruzada categórica; a regularização L2 foi definida como zero, o que desativa a penalização adicional de pesos.

Ao término de cada época, o código calcula perda e acurácia nos conjuntos de treino e validação sem executar atualizações nessas avaliações. A perda de validação é a métrica monitorada pelo \textit{early stopping}. A configuração permite quinze épocas sem melhora superior a 0,0001 antes da interrupção. Quando ocorre uma melhora suficiente, o programa copia todos os pesos e biases. Ao fim do treinamento, ele restaura a cópia da melhor época antes de usar o conjunto de teste. Assim, o teste não participa do ajuste nem da escolha da quantidade efetiva de épocas.

A etapa de avaliação utiliza as previsões do modelo restaurado e os rótulos de teste para calcular acurácia, precisão, recall, F1 por classe, relatório de classificação e matriz de confusão. As métricas são calculadas com funções auxiliares de scikit-learn, mas essas funções não treinam a rede. As curvas de perda e acurácia registram as métricas por época; a matriz de confusão apresenta as quantidades por classe real e prevista. Modelo, scaler, metadados de arquitetura, mapeamento de classes, configuração, métricas e gráficos são salvos em diretórios de resultados para permitir inspeção e reprodução.

\section{Resultados e discussão}

A execução principal documentada utilizou a seed 42, divisão estratificada de 70\% para treino, 15\% para validação e 15\% para teste, arquitetura com dezesseis entradas, duas camadas intermediárias com dezesseis e oito neurônios e sete saídas. Foram usadas ReLU nas camadas intermediárias, softmax na saída, lotes de 32 exemplos e taxa de aprendizado 0,001 com Adam. O limite configurado foi de cem épocas, mas a parada antecipada encerrou o treinamento após 59 épocas. Como a época 44 apresentou a menor perda de validação, seus parâmetros foram restaurados antes da avaliação final. A escolha da época, portanto, foi baseada na validação, sem utilizar o conjunto de teste para esse fim.

No teste dessa execução, composto por 2.042 amostras, a rede classificou corretamente 1.881 feijões e errou 161, alcançando acurácia de 92,12\%. A precisão média macro foi de 93,42\%, o recall macro foi de 93,21\% e o F1 macro foi de 93,29\%. O F1 ponderado foi de 92,15\%. A diferença entre as médias macro e ponderada está relacionada à quantidade desigual de exemplos por variedade: a média macro considera as classes com o mesmo peso, enquanto a ponderada leva em conta o suporte de cada uma.

A matriz de confusão também evidencia que o resultado geral não representa igualmente todas as variedades. As 78 amostras Bombay do teste foram classificadas corretamente, mas esse número relativamente pequeno recomenda cautela antes de generalizar esse desempenho. Horoz apresentou recall de 96,19\%, e Seker, de 94,08\%. A maior confusão ocorreu entre Dermason e Sira: 51 amostras Dermason foram previstas como Sira, e 26 amostras Sira foram previstas como Dermason. Sira teve precisão de 83,73\% e recall de 89,87\%. Isso indica que a rede reconheceu a maior parte dos exemplos Sira, mas também atribuiu essa classe a alguns feijões de outras variedades. A semelhança dos descritores geométricos pode contribuir para essas confusões, embora a matriz, sozinha, não determine sua causa física.

Além dessa execução principal, foram comparadas cinco configurações de arquitetura e taxa de aprendizado. A Tabela~\ref{tab:experimentos} apresenta a quantidade de camadas intermediárias, os neurônios de cada camada, a taxa de aprendizado e a acurácia obtida no teste. Os valores de acurácia foram convertidos para porcentagem para facilitar a leitura.

\begin{table}[ht]
\centering
\caption{Acurácia no teste para cinco configurações da MLP.}
\label{tab:experimentos}
\small
\setlength{\tabcolsep}{3pt}
\begin{tabular}{|p{1.25cm}|p{2.35cm}|p{3.25cm}|p{2.35cm}|p{3.15cm}|}
\hline
Execução & \shortstack{Número de\\camadas intermediárias} &
\shortstack{Neurônios por\\camada intermediária} &
\shortstack{Taxa de\\aprendizado} &
\shortstack{Acurácia\\no teste} \\
\hline
\rule{0pt}{2.5ex}1 & 1 & [16] & 0,001 & 92,21\% \\
\hline
\rule{0pt}{2.5ex}2 & 3 & [16, 16, 8] & 0,001 & 91,58\% \\
\hline
\rule{0pt}{2.5ex}3 & 2 & [16, 8] & 0,00001 & 82,32\% \\
\hline
\rule{0pt}{2.5ex}4 & 2 & [16, 8] & 0,1 & 89,03\% \\
\hline
\rule{0pt}{2.5ex}5 & 1 & [8] & 0,05 & 92,12\% \\
\hline
\end{tabular}
\end{table}

A primeira execução da tabela obteve a maior acurácia entre as cinco configurações, com 92,21\%, usando apenas uma camada intermediária de dezesseis neurônios. A segunda configuração, com três camadas intermediárias e a mesma taxa de aprendizado de 0,001, alcançou 91,58\%. Nesse conjunto de experimentos, acrescentar camadas não melhorou o resultado: a rede mais profunda ficou 0,63 ponto percentual abaixo da configuração de uma camada. Isso não significa que redes com mais camadas sempre tenham desempenho inferior; mostra que o aumento da profundidade, por si só, não garante uma classificação melhor. A arquitetura precisa ser avaliada com os dados e as demais escolhas de treinamento utilizados.

A taxa de aprendizado também alterou o comportamento observado. Na configuração com duas camadas intermediárias e taxa 0,00001, a acurácia final de teste foi 82,32\%, a menor da tabela. Esse valor ficou abaixo da média aritmética das cinco acurácias comparadas, que foi de 89,45\%. O histórico dessa execução mostra uma evolução lenta: a acurácia de validação começou em 5,24\% e chegou a 83,74\% ao fim das cem épocas. No registro por época, ela cresce gradualmente, sem oscilações contínuas. Portanto, os dados salvos sustentam que essa taxa muito baixa levou a uma aprendizagem lenta e a um resultado final inferior, mas não confirmam uma sequência de subidas e descidas persistentes na curva por época.

As taxas mais altas não produziram automaticamente a maior acurácia. Com taxa 0,05 e uma camada de oito neurônios, a quinta execução alcançou 92,12\%; durante o treinamento, a acurácia de validação permaneceu aproximadamente entre 90,40\% e 92,41\%, próxima da faixa de 91\% a 93\%. Já a taxa 0,1, na quarta execução, terminou com 89,03\% de acurácia no teste. Esse resultado mostra que elevar a taxa não garante aumento dos acertos. Também é necessário distinguir a estabilidade observada durante o treinamento do desempenho final no teste: são medidas relacionadas, mas não equivalentes. Como as arquiteturas e as taxas variam entre algumas linhas, os resultados devem ser interpretados como comparação das configurações completas, sem atribuir toda diferença a um único parâmetro isoladamente.

Em conjunto, os experimentos mostram que a qualidade da MLP depende da combinação entre arquitetura e parâmetros de treinamento. A acurácia informa a proporção total de acertos, enquanto precisão, recall, F1 e matriz de confusão ajudam a entender quais classes geram erros. Os resultados correspondem às partições e configurações registradas e não garantem o mesmo desempenho em novas amostras. Para tornar futuras comparações mais conclusivas, convém alterar um parâmetro por vez e manter constantes a divisão dos dados, a seed e os demais parâmetros, registrando também as curvas de treino e validação.

\section{Conclusão}

O projeto implementou e avaliou uma MLP para classificar variedades de feijão a partir de dezesseis atributos numéricos. A execução examinada classificou corretamente 1.881 das 2.042 amostras de teste e atingiu 92,12\% de acurácia. A análise por classe complementou essa medida e destacou a confusão entre Dermason e Sira, mostrando por que a inspeção da matriz de confusão é necessária para compreender onde o classificador falha. Esses valores descrevem a execução registrada e não devem ser interpretados como uma taxa garantida para dados futuros.

Do ponto de vista dos autores, a principal qualidade da implementação é tornar explícitas as etapas que bibliotecas de alto nível normalmente ocultam. Os métodos que processam uma camada, guardam os valores intermediários, calculam os gradientes e atualizam pesos e biases podem ser examinados diretamente no código. Essa transparência favorece a explicação dos conceitos de neurônio, peso, bias, camada, lote e época e permite modificar a arquitetura sem substituir o algoritmo central por uma caixa-preta.

O resultado também deve ser interpretado considerando suas limitações. A avaliação foi feita em uma única divisão estratificada e com uma arquitetura e uma configuração específicas; não houve busca sistemática de hiperparâmetros nem validação externa em outra coleta. O conjunto contém classes com quantidades diferentes de exemplos e os atributos disponíveis descrevem geometria, não cor, textura ou outras informações visuais. A MLP recebe descritores já extraídos e não substitui um processo de visão computacional. Estudos futuros podem preencher a tabela experimental com comparações controladas de arquiteturas, lotes e taxas de aprendizado, repetir a divisão com outras sementes e examinar intervalos de variação das métricas.

Em síntese, o experimento demonstra que uma rede implementada manualmente com NumPy consegue aprender uma relação entre as medidas geométricas e as variedades presentes no Dry Bean Dataset, mantendo um fluxo de avaliação separado entre treino, validação e teste. Os acertos observados são relevantes como evidência do funcionamento da implementação, enquanto os erros e as limitações indicam que a qualidade do modelo deve ser avaliada por classe e em experimentos adicionais. A tabela permanece sem valores para receber essas comparações quando forem executadas.

\bibliographystyle{sbc}
\bibliography{referencias}

\end{document}
